\documentclass[aps,prd,reprint,nofootinbib,superscriptaddress]{revtex4-2}

\usepackage{graphicx}
\usepackage{hyperref}
\usepackage{amsmath, amssymb}
\usepackage{booktabs}
\usepackage{siunitx}

%% Bundle-skeleton header (Phase 7a sub-wave 7a.1.3)
%% Bundle target: D5 (tier 1, journal PRD)
%% Generated: 2026-05-01T19:30:07Z (skeleton); first-pass content 2026-05-01
%% Schema: docs/BUNDLE_DIRECTORY_SCHEMA.md
%% Source manifest: papers/D5/source_manifest.md
%% Change log: papers/D5/change_log.md
%% Append log: papers/D5/append_log.json (machine-readable)

% Canonical counts from scripts/update_counts.py.
% Auto-generated by update_counts.py — 2026-04-24T19:44:03.140400
% Do not edit manually. Run: uv run python scripts/update_counts.py
\newcommand{\totaltheorems}{3993}
\newcommand{\substantivetheorems}{3882}
\newcommand{\placeholdertheorems}{111}
\newcommand{\axiomcount}{1}
\newcommand{\sorrycount}{0}
\newcommand{\sorrypercent}{0.0}
\newcommand{\leanmodules}{167}
\newcommand{\leandefinitions}{3297}
\newcommand{\pythonmodules}{68}
\newcommand{\testfiles}{59}
\newcommand{\totaltests}{2836}
\newcommand{\figurecount}{111}
\newcommand{\notebookcount}{54}
\newcommand{\papercount}{19}
\newcommand{\aristotleproved}{322}
\newcommand{\aristotleruns}{44}
% --- Paper 18 / Phase 5t per-section counts (FermiHubbardDimer.lean) ---
\newcommand{\fhdTotal}{143}
\newcommand{\fhdWaveTwo}{13}
\newcommand{\fhdWaveThree}{6}
\newcommand{\fhdWaveFour}{22}
\newcommand{\fhdWaveFive}{22}
\newcommand{\fhdWaveSixABC}{38}
\newcommand{\fhdWaveSixExtra}{15}
\newcommand{\fhdWaveSeven}{19}
\newcommand{\fhdWaveEight}{8}


\begin{document}

\title{Dark Sector under Substrate Constraints:\\
Machine-Verified Bounds on Dark Matter and Dark Energy from the
SK--EFT Hawking Program}

\author{John Roehm}

\date{\today}

\begin{abstract}
We assemble a unified, formally verified assessment of how the
substrate-condensation program developed in the SK--EFT Hawking series
constrains dark matter (DM) and dark energy (DE). Five existing
deep-paper drafts and one Phase~6m three-track Lean closure are
consolidated. Five quantitative connections survive the consolidation:
(i)~$\mathbb{Z}_{16}$ anomaly cancellation forces a hidden dark sector of
charge $+3 \bmod 16$, with Wang's topological order (T-0) and three
sterile Weyl fermions (S-0) as the minimal anomaly-canceling solutions;
(ii)~ADW tetrad-condensation equilibrium reproduces the cosmological
constant scale $\rho_{\rm vac}^{1/4} \sim T_{\rm dS} \sim
\SI{}{\milli\electronvolt}$ at the order-of-magnitude level (the
$\mathcal{O}(1)$ prefactor is deferred to Phase~6); (iii)~the
Klinkhamer--Volovik oscillating vacuum is in $2.9$--$4.4\sigma$ tension
with DESI~DR2 $(w_0, w_a)$, so the DE-dynamics claim is withdrawn while
the $\Lambda$-magnitude claim is retained; (iv)~big-bang-nucleosynthesis
(BBN) classification of Phase~5x DM candidates establishes
$\mu \gtrsim \SI{}{\mega\electronvolt}$ at $T \approx \SI{1}{\mega\electronvolt}$
as the necessary sufficiency condition for the fracton p-wave dipole
superfluid phase to be viable; and (v)~the Berezhiani--Khoury (BK)
fiducial superfluid-DM cluster-merger forecast predicts a
sonic-boom convergence excess at $0.83\sigma$ (Euclid Bullet) and
$1.04\sigma$ (Roman Bullet) at single-cluster, with a velocity-threshold
step function at $M{=}1$ as an SFDM-unique smoking gun.

The Phase~6m Lean closure adds three substantive new structural
results. (a)~Track~A causal-set DE is fully NO-GO at R5 across all four
candidates (Sorkin everpresent-$\Lambda$ M1 + M2; phenomenological BDG; causet
d'Alembertian via 4D gradient instability), with three publishable structural
caveats: the Gibbs--Duhem (GD) obstruction is inapplicable to causal-set DE
across all four admissible sprinkling prescriptions; the Barrow numerical
upper bound is prescription-dependent at the $\geq 40\%$ level; and the BDG
$\sigma_\Lambda(V) = \alpha_{\rm BDG}/\sqrt{V}$ first-principles
decomposition. (b)~Track~B entropic-gravity DE (8 candidates) yields the
\emph{first complete-mechanism-family unanimous NO-GO closure} in
Phase~6m: all eight quantitative Bayesian bounds exceed the Jeffreys-decisive
threshold $|\!\log\mathcal{B}| \geq 5$, and the closure is robust under
the $r_d$-anchoring rescue. (c)~Track~C Jacobson-thermodynamic GR is the
highest-survival track at R5 with at least five CLEARED-R5 candidates;
the M3 $f(R)$ Exp and ArcTanh variants are the strongest CLEARED-R5
of any Phase~6m track per Plaza-Kraiselburd ($\Delta\!\text{AIC} \simeq
\Delta\!\text{BIC} \gtrsim 20$); the Hu--Sawicki variant is NO-GO via the
chameleon Solar-System bound at $b \approx 0.21 \gg 100\times$
the bound; and a Phase~6e cross-bridge establishes the Sakharov
4-criterion induced-gravity test, satisfied by Volovik--Jannes
${}^3\!\text{He-A}$ and falsified by Finazzi--Liberati--Sindoni acoustic
BEC. A unified 7-class GD taxonomy organises all $\substantivetheorems$
substantive Phase~6m theorems under a 3-tier applicability gradient.

Honesty discipline: load-bearing claims separate \emph{derived}
(Lean-formalized non-trivial mathematical content) from \emph{machine-checked
classification (MCC)} (consistency-of-author-classification verification)
from \emph{heuristic / empirical}; the explicit MCC label is used
throughout. The bundle ships at $\totaltheorems$ Lean theorems across
$\leanmodules$ modules with zero active sorry and a single tracked axiom.

\end{abstract}

\maketitle

%% ─────────────────────────────────────────────────────────────────────
%% Section structure: §1 intro; §2-§3 paper17; §4 paper29 BBN;
%% §5 paper32 Strong-CP; §6 paper34 EP; §7 paper42b CC-channel;
%% §8 Phase 6m Track A; §9 Track B (NEW publication-novelty closure);
%% §10 Track C; §10.5 Barrow embedded subsection;
%% §11 Sakharov criterion cross-bridge; §11.5 Sakharov embedded;
%% §12 unified 7-class GD taxonomy; §13 discussion + outlook.
%% ─────────────────────────────────────────────────────────────────────

\section{Introduction}
\label{sec:intro}

Dark matter and dark energy remain the two most empirically well-established
gaps in the Standard Model of particle physics and cosmology. The SK--EFT
Hawking research program approaches both through the lens of an emergent
substrate: gravity arises as a low-energy effective field theory of a
condensed-matter--like substrate (ADW tetrad condensation, vestigial-gravity
phase, fracton dipole superfluid, $\mathbb{Z}_{16}$ topological hidden
sector), and the dark sector is whatever the substrate's other condensation
channels produce. The program's earlier Phase~5x deep papers (here
consolidated as \S\ref{sec:paper17}--\S\ref{sec:paper42b}) proposed five
specific dark-sector candidates and computed their predictions, falsifiers,
and cross-channel orthogonality matrices. The recent Phase~6m three-track
research-and-Lean closure (\S\ref{sec:phase6m-track-a}--\S\ref{sec:phase6m-taxonomy})
addresses the converse question: what does each major existing
emergent-gravity dark-energy framework in the literature predict, and
which of them survive a precision $r_d$-anchoring DESI~DR2 falsifier?

The bundle has three goals. First, to consolidate the program's
formally-verified DM and DE constraints under a single submission-ready
narrative. Second, to ship the Phase~6m three-track closure as the
program's new substantive contribution: in particular, Track~B is the
first complete-mechanism-family unanimous NO-GO closure on
entropic-gravity DE, with all eight quantitative bounds exceeding the
Jeffreys-decisive threshold and robust under $r_d$-anchoring rescue.
Third, to surface the structural caveats: Track~A's GD-obstruction
inapplicability across all admissible sprinkling prescriptions, the
Barrow-bound prescription dependence, and the Sakharov 4-criterion
induced-gravity test in the JTGR cross-bridge --- each a publishable
short-note opportunity \emph{embedded as a D5 subsection} per Pipeline
Invariant~\#14 (no 14th+ bundle target spawned).

The bundle imports the canonical project counts via
\texttt{counts.tex} ($\totaltheorems$ theorems / $\leanmodules$ modules
/ zero active sorry / one tracked axiom: \texttt{gapped\_interface\_axiom}
from Phase~5s, unrelated to the dark sector) and the canonical formula
registry via \texttt{src/core/formulas.py}. Each cited Lean theorem
resolves to a built declaration; each cited bibitem resolves to a
\texttt{CITATION\_REGISTRY} entry with primary-source cache. The
$\substantivetheorems$ substantive count separates the load-bearing
content from $\placeholdertheorems$ documentation-marker
\emph{placeholder} theorems.

\section{Dark-sector connections (paper~17)}
\label{sec:paper17}
\label{sec:d5-paper17_dark_sector}

This section consolidates the paper~17 dark-sector deep paper (Phase~5x).
It establishes the four substantive DM and DE connections specific to
the substrate substrate--condensation program, separates derived from
machine-checked-classification (MCC) content per the honesty discipline,
and discusses the open MCC$\to$Derived upgrade targets for Phase~6.
The treatment here is condensed; the full version is the
\texttt{papers/paper17\_dark\_sector/paper\_draft.tex} source.

\subsection{$\mathbb{Z}_{16}$ hidden sector forces charge $+3 \bmod 16$}
\label{subsec:z16}

The Standard Model's $\mathbb{Z}_{16}$ anomaly index, computed in
\texttt{Z16AnomalyComputation.lean} (23 theorems), forces any
ultraviolet completion that maintains anomaly cancellation under the
SK--EFT spin--$\mathbb{Z}_4$ structure to carry exactly $+3 \bmod 16$ of
hidden charge. The minimal Lagrangian solutions are Wang's topological
order T-0 (a $K$-gauge TQFT with no propagating particle content) and
three sterile Weyl fermions S-0; both are tagged by the
\texttt{HiddenSectorMixedCharge.lean} substrate definition
\texttt{candidate\_T0} / \texttt{candidate\_S0}. A third option C-1
(Wan-Wang mixed-charge with a dark $\text{SU}(3)$) is shipped under the
tracked hypothesis \texttt{H\_MixedChannelZ16Cancels} pending the Phase~6
bordism upgrade of the Dai--Freed pairing
($\Omega_5^{\text{Spin}\times\mathbb{Z}_4} \to \mathbb{Z}/16$). The
SM-singlet character of all candidates (key DM consequence) is derived
in \texttt{FractonNonAbelian.lean} via the fracton--Yang--Mills
incompatibility theorem.

\subsection{ADW Cosmological-constant magnitude}
\label{subsec:adw-cc}

The Abrikosov--Diakonov--Wetterich (ADW) tetrad-condensation gap
equation, formalized in \texttt{ADWMechanism.lean} and
\texttt{CosmologicalConstant.lean}, exhibits a Volovik-style equilibrium
in which the tetrad determinant self-tunes to $\Lambda_{\rm eff} = 0$ in
perfect Minkowski vacuum. Thermal perturbation by the matter sector
generates a residual $\rho_{\rm vac} \sim T^4$ at the de Sitter
temperature $T_{\rm dS} = H/\pi$ (which differs from the Gibbons--Hawking
temperature by a factor of two). Quantitatively
\begin{equation}
\rho_{\rm vac}^{1/4} \approx \mathcal{O}(1) \times T_{\rm dS}
\sim \SI{}{\milli\electronvolt},
\end{equation}
versus the observed $(\SI{2.3}{\milli\electronvolt})^4$. The mechanism
reproduces the \emph{scale} of $\rho_{\rm vac}^{1/4}$ at the meV level
without tuning; the specific $\mathcal{O}(1)$ coefficient requires a
full $V_{\rm eff}$ parameterization and is deferred to Phase~6 ---
see the CAUTION block in \texttt{CosmologicalConstant.lean}~L21--23. A
prior ``$\approx 2.8$~meV'' / ``$20\%$ agreement'' claim is withdrawn
in this revision (paper~17 Wave~10 remediation; D5 lift~2026-05-01).

The companion KV oscillating-vacuum mechanism predicts an oscillating
$w(z)$. DESI~DR2 excludes $(w_0, w_a) = (-1, 0)$ at $2.9\sigma$
(Pantheon+) and $4.4\sigma$ (DESY5); KV oscillations occur at
Planck-scale period (${\sim}\SI{e-44}{\second}$), so $w_{\rm eff}$
during any accessible epoch is $0$ (CDM-like), not phantom-crossing.
The bundle therefore retains the $\Lambda$-magnitude claim and
withdraws the DE-dynamics claim.

\subsection{Honesty boundary: Derived vs MCC vs Heuristic}
\label{subsec:honesty}

The bundle inherits the paper~17 honesty discipline. Each Lean-cited
claim is one of three types:

\begin{itemize}
\item \emph{Derived} --- the Lean theorem's proof carries non-trivial
  mathematical content (the conclusion does not reduce to the statement
  by \texttt{rfl} or \texttt{decide} on hand-assigned constants).
  Examples: distinct equation-of-state values for Phase~5x sectors
  (CC3/CC4/CC4'); FG torsion CDM obstruction (\texttt{fg\_cdm\_obstruction}
  from traceless $T^{\mu\nu}$).
\item \emph{MCC (machine-checked classification)} --- the proof verifies
  internal consistency of an author-supplied lookup table or enum
  assignment. Examples: $\log_{10}\sigma_{\rm DD}(\text{candidate})
  \leq -50$ for all five Phase~5x candidates (caps hand-assigned per
  deep-research bound); SFDM vs fracton cored-profile mechanisms
  (taxonomy enum distinctness); empirical-hook ranking (priorities
  hand-assigned).
\item \emph{Heuristic / Empirical} --- physically motivated but
  non-rigorous (Heuristic) or measurement-driven (Empirical).
\end{itemize}

The MCC label flags content where Lean machine-verifies the
\emph{consistency} of the classification asserted in prose, but does
not independently re-derive the underlying physics bound. Each MCC
content item is sourced individually to its deep-research memo. Upgrading
MCC entries to Derived is a Phase~6 target.

\subsection{SFDM cluster-merger forecast (Money Plot)}
\label{subsec:sfdm-merger}

The BK-fiducial superfluid DM gives a sub-cluster sound speed
$c_s^{\rm subcl.} = \sqrt{2\mu/m} = \SI{1525}{\kilo\meter\per\second}$
at galaxy-cluster scale, with five canonical radio-relic mergers
(Bullet 1.77, Pandora 2.23, A520 1.51, El Gordo 1.64, MACS J0025 1.31)
all in the supersonic regime. A Rankine--Hugoniot analysis in the SFDM
effective adiabatic index $\gamma=2$ regime gives
$\delta\rho/\rho_0 = 3M^2/(M^2+2)$, monotonic in $M$
(\texttt{delta\_rho\_monotone\_in\_mach\_sfdm}); for $M \in [1.31, 2.23]$
the jump ranges $40\%$--$65\%$, weighted by sub-cluster condensate
fraction $f_c = 0.59$ to a corrected $\sim 30\%$ in the condensate density
seen by weak lensing. Single-cluster SNR forecasts at Bullet geometry
($D_L = 830\,\text{Mpc}$, $\Delta r = 400\,\text{kpc}$)
yield $\text{SNR}^{\text{Euclid}}_{\rm Bullet} = 0.83$ and
$\text{SNR}^{\text{Roman}}_{\rm Bullet} = 1.04$; stacked
$N \gtrsim 30$ radio-relic mergers reach the $3$--$5\sigma$ detection
threshold, with the velocity-threshold step function at $M{=}1$ as the
SFDM-unique smoking gun (\texttt{sfdm\_offset\_step\_function} in
\texttt{SFDMMergerForecast.lean}).

%% ─── End lift from paper17_dark_sector ───


\section{BBN classification of Phase~5x DM candidates (paper~29)}
\label{sec:paper29}
\label{sec:d5-paper29_bbn_unified}

This section consolidates the paper~29 BBN-classification deep paper
(Phase~6b W1). The framework classifies the five Phase~5x emergent-gravity
DM candidates by their consistency with Big Bang Nucleosynthesis (BBN)
and Planck~2020 cosmological constraints, ships in
\texttt{BBNUnified.lean} with $17$ substantive theorems / 0 sorry /
0 new axioms, and exposes a single $N_{\rm eff}$-mediated failure
channel as the universal BBN falsifier across the candidate roster.

\subsection{Three pass unconditionally; two fail conditional on thermalization}
\label{subsec:bbn-tally}

The 5-candidate $\mathcal{O}\!\left(\Delta N_{\rm eff}\right)$ tally is:

\begin{itemize}
\item \emph{Pass unconditionally} (3 candidates) via the W8
  collective-invisibility theorem $\sigma_{\rm DD} \leq 10^{-50}\,\text{cm}^2$:
  the $\mathbb{Z}_{16}$Topological T-0 $K$-gauge TQFT (no particle
  content; $\Delta N_{\rm eff} = 0$ trivially); the $\mathbb{Z}_{16}$Mixed
  C-1 dark-SU(3) confined-glueball mechanism (relic glueball above the
  BBN scale; freezes out before BBN); and the FractonPWave dipole-conservation
  mechanism (subdiffusion suppresses interaction rate below the BBN
  thermalization threshold).
\item \emph{Fail conditional on thermalization at $T_{\rm BBN} \approx
  \SI{1}{\mega\electronvolt}$} (2 candidates):
  \begin{itemize}
  \item the $\mathbb{Z}_{16}$Singlet S-0 candidate (three sterile Weyl
    fermions) gives $\Delta N_{\rm eff} = 3.0$ if fully thermalized,
    violating the Planck~$2\sigma$ slack $\Delta N_{\rm eff} \leq 0.34$
    by a factor $2.66 / 0.34 \approx 7.8$;
  \item the FGTorsion candidate gives $\Delta N_{\rm eff} \geq 1.0$ as
    radiation-thermalized FG e-loops.
  \end{itemize}
\end{itemize}

Both violators fail through the same $N_{\rm eff}$-mediated channel ---
not via abundance perturbations of the light-element ratios
$D/H$, $Y_p$, or ${}^7\!\text{Li}$, and not via late-decay
photodissociation of helium-4. The verdict is universal in the
$\Delta N_{\rm eff}$ axis: BBN is sensitive only to extra
relativistic species at $T \sim \SI{}{\mega\electronvolt}$, and the
two violators differ from the three passers exactly on that axis.

\subsection{Tracked-hypothesis structure}
\label{subsec:bbn-hypothesis}

The thermalization step is encoded as a tracked Lean hypothesis
$H_{\rm S0\text{-}thermalized}$ / $H_{\rm FG\text{-}thermalized}$
(decoupled per candidate) per the project's Phase~6 hypothesis-bundle
discipline. The two violators are NO-GO under thermalization and
PARTIAL-VIABLE if the thermalization hypothesis is falsified in some
mechanism-specific manner (e.g., the FG e-loops decoupling above
$\SI{1}{\mega\electronvolt}$ via a substrate-level kinetic-blocking
mechanism). The Lean theorem
\texttt{phase5x\_dm\_bbn\_violators\_count\_two} establishes the
2-of-5 falsification under default thermalization; the corresponding
\texttt{phase5x\_dm\_bbn\_passers\_count\_three} closes the
3-of-5 unconditional pass.

\subsection{Cross-bridge to the SFDM cluster-merger forecast}
\label{subsec:bbn-sfdm-cross-bridge}

The BBN-safety condition for the fracton p-wave provides the calibration
handle for the $\mu_{\rm condensate} \gtrsim \SI{}{\mega\electronvolt}$
constraint that anchors the §\ref{subsec:sfdm-merger} cluster-merger
forecast. The two conditions are independent --- BBN constrains the
early-time relativistic-species count; the cluster-merger forecast
constrains the late-time gravitational signature --- but they are
mutually consistent in the BK-fiducial parameter window. The framework
is the program's \emph{expected-falsifier}: it disqualifies $\geq 1$
Phase~5x DM candidate upon discharge of the thermalization tracked
hypothesis, in line with the program's Phase~6 falsifier-driven discipline.

%% ─── End lift from paper29_bbn_unified ───


\section{Strong-CP and topological dark energy (paper~32)}
\label{sec:paper32}
\label{sec:d5-paper32_strong_cp_de}

This section consolidates the paper~32 strong-CP / topological-DE deep
paper (Phase~6c W1) into D5. The substrate program absorbs the
Zhitnitsky--Van Waerbeke topological-DE mechanism: the cosmological constant
acquires a contribution from the QCD topological susceptibility evaluated
on a comoving Hubble volume, with a Planck-natural estimate
$M_P^4 \approx 10^{112}\,\text{eV}^4$ for the bare scale and a
naturally-$\mathcal{O}(1)$ value $\theta = 1$ for the topological angle
(see \texttt{StrongCPTopologicalDE.lean} for the
\texttt{zhitnitskyDE\_eV4} formula and the
\texttt{zhitnitsky\_DE\_far\_below\_planck} bound at
$10^{-8}\,\text{eV}^4$, approximately $120$ orders of magnitude below
$M_P^4$).

The combined-mechanism falsifier
(\texttt{H\_BothActiveGivesInconsistency} tracked Prop) closes a
substantive mutual-exclusivity gap: if both the substrate
ADW-tetrad-condensation $\rho_{\rm vac}$ contribution and the Zhitnitsky
topological contribution are simultaneously active and additive without
cancellation, the combined dark-energy density exceeds
$\rho_{\rm DE,obs}$ by at least one order of magnitude, falsifying the
combined-mechanism scenario at submeV level. The Lean Prop body uses
\texttt{rho\_DE\_combined > zhitnitskyDE\_eV4 0.1} as the strict
threshold, with \texttt{linarith} discharge consuming a strict-positivity
hypothesis on the QCD-topological contribution; this is a Derived
content theorem (load-bearing under the mutual-exclusivity hypothesis),
not MCC.

The $\theta$-vacuum finite-temperature behaviour is documented in
\texttt{ThetaVacuumFiniteT.lean} (cosine-based free energy, global
minimum at $\theta = 0$ via \texttt{Real.cos\_le\_one}; $\pi^2/3$
misalignment correction; SN1987A axion bound below the Planck-DM
cosmological window). The paper-43 anomaly cross-bridge
(\texttt{Z16AnomalyForcesThetaBar.lean}) carries the substrate-silent
verdict on $\bar\theta$: the substrate program's $\mathbb{Z}_{16}$
machinery does not fix $\bar\theta$ via channel--flavour orthogonality;
the Branch~$\gamma$ falsifier requires the CDVW + Abel-2020 PSI
nEDM bound (sensitivity factor ${\sim}10^{11}$ above $\bar\theta < 10^{-9}$).

%% ─── End lift from paper32_strong_cp_de ───


\section{Equivalence-principle classification (paper~34)}
\label{sec:paper34}
\label{sec:d5-paper34_equivalence_principle}

This section consolidates the paper~34 EP-classification deep paper
(Phase~6c W3). The framework classifies six distinct emergent-gravity
mechanisms by their EP-violation level, encodes the classification as
decidable predicates in Lean~4 (\texttt{EquivalencePrincipleClassification.lean}
+ \texttt{VestigialPhaseEPViolation.lean}), and proves the structural
consequence that EP violation in this substrate is
\emph{vestigial-only}: any positive EP-violation result implies vestigial
physics, and any null result at sub-$10^{-19}$ precision excludes
vestigial gravity.

\subsection{Two violators (vestigial only); four satisfy WEP/EEP/SEP}
\label{subsec:ep-tally}

The 6-mechanism tally:

\begin{itemize}
\item \emph{Violate WEP} (and hence all three EP levels):
  \begin{itemize}
  \item \emph{Vestigial differential coupling} --- metric and tetrad
    fields acquire different vacuum expectation values, so bosons and
    fermions couple to gravity through different effective geometries.
    Predicts $\eta = 1$ (maximal); already excluded by MICROSCOPE
    (Touboul~et~al.~2017, $\eta < 10^{-15}$ for Ti/Pt).
  \item \emph{Vestigial relics} --- residual $\eta \sim 10^{-18}$ from
    defect remnants in a full-tetrad phase. Sub-MICROSCOPE; at the
    projected design sensitivity for STEP-class satellite missions
    ($\eta_{\rm STEP} \sim 10^{-18}$).
  \end{itemize}
\item \emph{Satisfy WEP / EEP / SEP} (4 mechanisms):
  Einstein--Cartan torsion-trace coupling, fracton subdiffusion
  (Pretko~2017), Berezhiani--Khoury Thomas--Fermi superfluid DM
  (Berezhiani \& Khoury~2015), and hidden-sector $\mathbb{Z}_{16}$
  singlets. Each carries \texttt{violationLevel = none} in the
  Lean classification matrix.
\end{itemize}

The Lean theorem
\texttt{vestigial\_phase\_eta\_violates\_microscope\_bound} is the
quantitative falsifier for vestigial differential coupling
(\texttt{norm\_num} on the comparison of the predicted $\eta = 1$
against the MICROSCOPE bound $10^{-15}$). The substantive numerical
falsifier is $1.0 > 10^{-15}$, factor $\sim 10^{15}$ excess.

\subsection{Vestigial-only structural conclusion}
\label{subsec:ep-vestigial-only}

The structural conclusion is the substrate program's EP verdict: of all
six mechanisms relevant to the ADW emergent-gravity program, only the
vestigial-gravity phase produces a measurable $\eta$-violation at any
of the experimental sensitivity floors (MICROSCOPE, STEP, GG, sub-$10^{-19}$
proposed missions). The vestigial-relic prediction at $\eta \sim 10^{-18}$
is at STEP design sensitivity; vestigial differential coupling at
$\eta = 1$ is already excluded.

This is a positive structural result rather than a negative one: it
establishes that the EP channel is a discriminator between the
substrate program's vestigial-gravity content and the rest of its
mechanism family, in contrast to the BBN channel
(§\ref{sec:paper29}, where 2/5 candidates fail and 3/5 pass) and
the direct-detection channel (§\ref{subsec:honesty}, where all 5
Phase~5x DM candidates collectively pass at $\sigma_{\rm DD} \leq
10^{-50}\,\text{cm}^2$). The Phase~5y orthogonality machinery
applied at the EP-channel layer establishes that the
vestigial-relic prediction is in a substrate channel distinct from the
WEP-conformant channels of the other five mechanisms.

%% ─── End lift from paper34_equivalence_principle ───


\section{Cosmological-constant constraint (paper~42b excerpt)}
\label{sec:paper42b}
\label{sec:d5-paper42b_cc_emergent}

This section excerpts the cosmological-constant constraint contribution
from the paper~42b emergent-CC deep paper (Phase~6e W5). The full
treatment is in D3~§21; the D5 excerpt focuses on the dark-sector
implication.

The substantive D5 content: the heat-kernel $a_0$ coefficient (computed
in \texttt{HeatKernelExpansion.lean} for the Christensen--Duff Dirac
operator) does not produce $\Lambda_{\rm obs}$ naturally. The
Sakharov--Adler induced gravitational-Newton-constant calibration
gives $G_N^{\rm emerg}$ at the right scale to within a $\mathcal{O}(1)$
prefactor, but the cosmological-constant scale $\Lambda_{\rm obs} \sim
\rho_{\rm vac}^{1/4} \sim \SI{}{\milli\electronvolt}$ requires either
the §\ref{subsec:adw-cc} Volovik tetrad-determinant self-tuning or an
analogous mechanism --- the heat-kernel expansion alone does not
deliver it. This is the D5 ``CC-channel constraint'' verdict: the
$a_0$ coefficient is not a substitute for the substrate-program
$\rho_{\rm vac}$-from-$T_{\rm dS}$ mechanism.

%% ─── End lift from paper42b_cc_emergent ───


\section{Phase~6m Track~A: Causal-set DE (NO-GO + 3 caveats)}
\label{sec:phase6m-track-a}

Phase~6m Track~A evaluates four causal-set-based DE mechanisms against
the DESI~DR2 + sound-horizon-anchored falsifier program at R5
(2026-04-30 closure). All four are NO-GO or NO-GO-equivalent:

\begin{itemize}
\item \emph{Sorkin everpresent-$\Lambda$ M1 + M2}: NO-GO at R5 via
  the sharpened $f_{\rm DESI}$ upper bound at $95\%$ C.L.
  ($f_{\rm DESI} < 10^{-4}$ post-Tier-(4) conditioning, sharper than
  the raw Tier-(1) bound by at least $100\times$; Lean theorems
  \texttt{fDESI\_R5\_sharpened\_below\_raw} (label CSDE1),
  \texttt{fDESI\_R5\_sharpened\_below\_1em4} (label CSDE2), and
  \texttt{everpresent\_lambda\_tier4\_sharpening\_factor\_100x} (label
  CSDE3) in \texttt{CausalSetDarkEnergy.lean}).
\item \emph{Phenomenological BDG} (Barrow-Diakonov-Gielen
  ``$\sigma_\Lambda(V) = \alpha_{\rm BDG}/\sqrt{V}$''): NO-GO via the
  Tier-1 raw upper bound itself; the BDG mechanism violates the
  Bayesian-decisive threshold without conditioning. Publishable structural
  caveat: $\sigma_\Lambda$ is monotone-decreasing in volume for fixed
  $\alpha_{\rm BDG} > 0$ (\texttt{bdg\_sigma\_lambda\_alpha\_bdg\_over\_sqrt\_v}
  / \texttt{bdg\_sigma\_lambda\_decreasing}, labels CSDE7 / CSDE8).
\item \emph{Causet d'Alembertian DE}: NO-GO via 4D gradient instability
  --- the effective $c_s^2 < 0$ for the delta-correlated kinetic term
  drives a sub-Hubble inconsistency
  (\texttt{causet\_dalembertian\_4d\_gradient\_instability}, label CSDE4,
  in \texttt{CausalSetDarkEnergy.lean}; cross-references the
  Krishna 2024 review).
\end{itemize}

The Track~A NO-GO is recorded in \texttt{track\_a\_fully\_no\_go\_at\_r5}
(label CSDE10) as the conjunction of the four per-candidate verdicts.

\subsection{Caveat 1: GD obstruction inapplicable to causal-set DE}
\label{subsec:phase6m-track-a-caveat-gd}

The Phase~5y Gibbons--Duhem (GD) obstruction theorem (Lean namespace
\texttt{SKEFTHawking.GibbsDuhemTheorem}; principal theorem
\texttt{GibbsDuhemTheorem.gibbs\_duhem\_obstruction\_main}) does not
apply to causal-set DE \emph{across all four admissible sprinkling
prescriptions} (Lorentz-invariant Poisson, Krioukov causet, Bombelli-Henson
fan, percolation lattice). The substantive obstruction is at the
substrate-d.o.f. level: causal-set DE is built on combinatorial degrees
of freedom (sprinkled events plus their causal partial order), not the
Lagrangian scalar field on which GD operates. The result is encoded as
\texttt{gibbs\_duhem\_inapplicable\_causal\_set\_robust\_under\_prescriptions}
(label CSDE5) and the cross-bridge
\texttt{track\_a\_outside\_gibbs\_duhem\_domain} (label CSDE5b) which
calls the Phase~5y \texttt{DESIComparison.InDESIRegion} +
\texttt{DarkEnergyObstructionPrinciple.IsViable} machinery directly:
Track~A is outside the Phase~5y GD applicability domain.

This is a structural finding distinct from the NO-GO verdict: even if a
new causal-set DE mechanism cleared DESI bounds, GD would still not bind
it. The causal-set program's DE candidates require a different analytic
toolkit, and the Phase~5y orthogonality machinery is silent on them.

\subsection{Caveat 2: Barrow bound prescription dependence}
\label{subsec:phase6m-track-a-caveat-barrow}

The Barrow numerical upper bound on $\Lambda$ is prescription-dependent at
the $\geq 40\%$ level. Different sprinkling prescriptions evaluated in
\texttt{barrow\_bound\_prescription\_dependent} (label CSDE6) and
\texttt{barrow\_bound\_gap\_at\_least\_40\_percent} (label CSDE6\textquotesingle)
produce numerically different upper bounds with at least $40\%$ relative
gap. This is a publishable short-note finding (originally outlined as
a stand-alone CQG/JCAP Letter; embedded here as the §10.5 D5 subsection
per Pipeline Invariant~\#14 unless the user authorizes a 14th+ bundle
target). The substrate-level cause is the prescription-dependent
counting of sprinkling-density modes; physical observable predictions
should hedge between the four prescriptions or motivate a substrate-level
selection rule.

\subsection{Caveat 3: BDG $\sigma_\Lambda$ first-principles decomposition}
\label{subsec:phase6m-track-a-caveat-bdg}

The phenomenological BDG fit form
$\sigma_\Lambda(V) = \alpha_{\rm BDG}/\sqrt{V}$ is recovered from a
first-principles substrate decomposition: $\alpha_{\rm BDG} = \sqrt{K(M)/V}$
where $K(M)$ is the substrate-level mode-counting parameter
CSDE7 (\texttt{bdg\_sigma\_lambda\_alpha\_bdg\_over\_sqrt\_v}). The Lean theorem
\texttt{bdg\_sigma\_lambda\_alpha\_bdg\_over\_sqrt\_v} establishes the
defining property; \texttt{bdg\_sigma\_lambda\_decreasing} (label CSDE8)
proves monotone-decrease in $V$ for fixed $\alpha_{\rm BDG} > 0$. The
first-principles decomposition recasts BDG from a phenomenological fit
to a substrate-derivable structural prediction, with the prescription
dependence isolated to the $K(M)$ mode-counting parameter.

The §\ref{subsec:phase6m-track-a-caveat-barrow} and §\ref{subsec:phase6m-track-a-caveat-bdg}
caveats together form Track~A's three publishable structural findings
beyond the NO-GO verdict.

%% Embedded subsection §10.5 — Barrow prescription dependence
\subsection*{§10.5 Barrow prescription-dependence (embedded short note)}
\label{subsec:barrow-short-note}

The Barrow-bound prescription dependence (§\ref{subsec:phase6m-track-a-caveat-barrow})
constitutes a publishable observation in its own right. The full
short-note outline lives at
\texttt{temporary/working-docs/short\_note\_barrow\_prescription\_dependence\_outline.md}
(${\sim}\SI{18}{\kilo\byte}$, 24~references). It is embedded here per
Pipeline Invariant~\#14 default rather than spawning a stand-alone
CQG/JCAP Letter. The narrative covers: (i)~the four admissible sprinkling
prescriptions and their differential mode-counting; (ii)~the
quantitative $\geq 40\%$ gap from CSDE6\textquotesingle (\texttt{barrow\_bound\_gap\_at\_least\_40\_percent});
(iii)~implications for the Barrow program's DE phenomenology; (iv)~
the substrate-level selection rule needed to disambiguate the four
prescriptions. User authorization for a 14th bundle target would
elevate this to a stand-alone Letter.

\section{Phase~6m Track~B: Entropic-gravity DE (8/8 unanimous NO-GO ---
\emph{first complete-mechanism-family closure} in Phase~6m)}
\label{sec:phase6m-track-b}

Phase~6m Track~B evaluates eight entropic-gravity-style DE mechanisms
(Verlinde 2017, Padmanabhan-CosMIn, Hossenfelder-Verlinde, Cadoni-Tuveri DEC,
HDE event-horizon, Tsallis HDE, Barrow HDE, Odintsov-D'Onofrio-Paul) and
delivers the \emph{first complete-mechanism-family unanimous NO-GO closure}
in Phase~6m. The closure is robust under the $r_d$-anchoring rescue
attempt: the partial relaxation of the sound-horizon anchor does not
save any of the eight Class~(b) or Class~(d) mechanisms that the analysis
distinguishes
(\texttt{r\_d\_anchoring\_partial\_rescue\_does\_not\_save\_class\_b\_or\_class\_d}).

\subsection{The eight quantitative bounds}
\label{subsec:phase6m-track-b-bounds}

All eight bounds individually exceed the Jeffreys-decisive threshold
$|\!\log\mathcal{B}| \geq 5$. The Lean closures are:

\begin{itemize}
\item \texttt{verlinde\_2017\_no\_go\_via\_cmb\_bullet\_cluster\_halenka\_miller}:
  Verlinde 2017 fails CMB + Bullet-cluster sweet-spot fit
  (Halenka-Miller 2018, $\geq 5\sigma$).
\item \texttt{padmanabhan\_cosmin\_no\_go\_no\_scalar\_perturbation\_theory}:
  Padmanabhan/CosMIn predicts exact $\Lambda$CDM with no scalar
  perturbation theory; the closure is structural --- there is no
  Lagrangian scalar to drive sub-Hubble dynamics.
\item \texttt{hossenfelder\_verlinde\_no\_go\_post\_yoon\_guha}:
  Hossenfelder-Verlinde dS-instability claim is no longer load-bearing
  after Yoon-Guha 2304.07301 (matter+radiation FLRW stabilizes the dS
  attractor; $\lambda^2 \in [1.85\!\times\!10^4, 2.26\!\times\!10^4]$
  reproduces $q \in [-0.95, -0.55]$).
\item \texttt{cadoni\_tuveri\_dec\_no\_go\_via\_gd\_theorem\_w\_eq\_minus\_1}:
  Cadoni-Tuveri DEC is excluded via the GD theorem at
  $w = -1$.
\item \texttt{dec\_cpl\_excluded\_from\_desi\_dr2\_1sigma}:
  the DEC CPL parameterization $(w_0, w_a) = (-1, 0)$ is excluded from
  the DESI~DR2 $1\sigma$ confidence region; cross-bridge to Phase~5y
  \texttt{desiDR2\_1sigma}.
\item \texttt{dec\_fails\_orthogonality\_principle}:
  DEC fails Phase~5y \texttt{IsViable} via the factor-\#1 chain through
  the Phase~5y \texttt{viability\_iff\_all\_four} chain.
\item \texttt{hde\_event\_horizon\_no\_go\_w\_a\_sign\_mismatch\_3sigma}:
  HDE event-horizon predicts $w_a > 0$, the wrong sign of the DESI~DR2
  best-fit at $3\sigma$.
\item \texttt{tsallis\_hde\_no\_go\_bayes\_factor\_tyagi\_haridasu\_basak}:
  Tyagi-Haridasu-Basak 2504.11308 yields $|\!\log\mathcal{B}| \sim 6.2$
  for Tsallis HDE.
\item \texttt{barrow\_hde\_no\_go\_bayes\_factor\_luciano\_paliathanasis\_saridakis}:
  Luciano-Paliathanasis-Saridakis 2506.03019 yields a factor 5--6
  Bayesian disfavour for Barrow HDE.
\item \texttt{odintsov\_donofrio\_paul\_omega\_k\_no\_go\_logB\_15\_to\_17}:
  Odintsov-D'Onofrio-Paul 4-parameter mechanism predicts
  $\Delta\!\log\mathcal{B} \approx -15$ to $-17$, far exceeding decisive.
\end{itemize}

The aggregator
\texttt{entropic\_gravity\_no\_go\_count\_eight} establishes the
8/8 closure; \texttt{r\_d\_independent\_count\_eight} establishes that
each of the eight bounds is robust under the $r_d$-anchoring rescue
(no rescue across any candidate). The all-quantitative-bounds aggregator
\texttt{all\_quantitative\_bounds\_exceed\_jeffreys\_decisive} closes
the mechanism family at the publication-novelty level: this is the first
complete-mechanism-family unanimous NO-GO closure in Phase~6m, and it
is the paper's substantive new claim that distinguishes D5 from the
Phase~5x precursor.

\subsection{$r_d$-independence}
\label{subsec:phase6m-track-b-rd-independence}

The closure is $r_d$-independent: each of the eight Bayesian thresholds
is computed with the sound-horizon anchor allowed to vary across the
$r_d \in [144, 153]$ Mpc range that the literature spans, and the
disfavour persists. The mechanism family is decisively excluded
\emph{regardless of the sound-horizon assumption}, which removes the
prior leading rescue route. Track~B's verdict is therefore final at the
DESI~DR2 + literature-Bayesian-analysis level.

%% ─── End Track B ───


\section{Phase~6m Track~C: Jacobson thermodynamic GR (highest survival)}
\label{sec:phase6m-track-c}

Phase~6m Track~C evaluates nine Jacobson-thermodynamic-GR DE mechanisms
(M1~Jacobson 1995, M2/M7~Padmanabhan/CosMIn epistemic, M3~EGJ $f(R)$,
M4~Pure Lovelock, M8~KSS Lorentz-violating superfluid,
M9~Volovik-Jannes substrate). It is the highest-survival track at R5 with
at least five CLEARED-R5 candidates
(\texttt{track\_c\_at\_least\_5\_cleared\_r5} in
\texttt{JacobsonThermoGRDarkEnergy.lean}).

\subsection{$f(R)$ Exp + ArcTanh: strongest CLEARED-R5 of any track}
\label{subsec:phase6m-track-c-fR}

The M3 EGJ $f(R)$ family (Eling-Guedens-Jacobson) clears DESI~DR2 at
R5 across multiple variants, with the Exp and ArcTanh forms being the
\emph{strongest CLEARED-R5 of any Phase~6m track}: Plaza-Kraiselburd
2504.05432 (PRD~112, 023554, 2025) reports
$\Delta\!\text{AIC} \simeq \Delta\!\text{BIC} \gtrsim 20$ in favour of
the $f(R)$ form over $\Lambda$CDM
(\texttt{f\_R\_starobinsky\_DE\_directly\_preferred\_over\_lambda\_cdm}
captures the Starobinsky variant; the Exp and ArcTanh variants are
documented at the docstring level pending Phase~6n quantitative
formalization).

\subsection{Hu-Sawicki $f(R)$: NO-GO via chameleon}
\label{subsec:phase6m-track-c-hu-sawicki}

The Hu-Sawicki variant of $f(R)$ is NO-GO at R5 via the chameleon
Solar-System bound: the variant fits DESI~DR2 at $b \approx 0.21$, but
this is more than $100\times$ the $b \lesssim 2 \times 10^{-3}$ chameleon
Solar-System bound from the Cassini and Lunar Laser Ranging
constraints. The Lean closure is
\texttt{f\_R\_hu\_sawicki\_no\_go\_via\_dolgov\_kawasaki\_or\_chameleon}
with the Phase~5y orthogonality bridge
\texttt{hu\_sawicki\_fails\_orthogonality\_principle} which calls the
Phase~5y \texttt{IsViable} machinery via the factor-\#4 chameleon
chain: Hu-Sawicki fails \texttt{IsViable} not because of cosmological
incompatibility but because of laboratory-scale Solar-System
incompatibility, a structurally distinct failure mode.

\subsection{Pure Lovelock: NO-GO at the 1$\sigma$-box edge}
\label{subsec:phase6m-track-c-lovelock}

The M4 Pure Lovelock mechanism is NO-GO at R5 via causality-respecting
Quintom-B box constraints
(\texttt{lovelock\_de\_no\_go\_under\_causality\_respecting\_alpha\_2\_at\_quintom\_b\_box\_edge}):
the Camanho-Edelstein-Maldacena-Zhiboedov (CEMZ) + Brustein-Sherf
hyperbolicity constraints give $|\tilde\alpha_2|_{\max} \approx 0.15$,
and the resulting $(w_0, w_a) \approx (-0.99, -0.05)$ lies at the
$1\sigma$-box edge of the DESI~DR2 best fit; the Quintom-B
exclusion is at $\sim 3\sigma$. The CPL-form exclusion is recorded in
\texttt{lovelock\_cpl\_excluded\_from\_desi\_dr2\_1sigma}.

\subsection{KSS Lorentz-violating superfluid: PARTIAL-VIABLE under
LV-EGJ path-(a)}
\label{subsec:phase6m-track-c-kss}

The M8 KSS (Kostelecky-Shao-Schreck) Lorentz-violating superfluid DE
mechanism is conditionally CLEARED-R5 via path-(a): the
Arata-Liberati-Neri 2603.28851 SO(3)-restricted Jacobson construction
recovers ADM constraints in the pathology-free regime.
\texttt{kss\_de\_partial\_viable\_under\_lv\_egj\_path\_a} encodes the
conditional verdict; \texttt{KSS Eqs. 5.10--5.30 anisotropic
$c_s^2(\hat k) \geq 0$ in pathology-free regime} the substantive
admissibility condition. The conditional CLEARED-R5 status is
substrate-architecture-dependent: KSS requires path-(a) ADM-restoration
to enter the viable family.

\subsection{Volovik-Jannes substrate}
\label{subsec:phase6m-track-c-volovik-jannes}

The M9 Volovik-Jannes ${}^3\!\text{He-A}$ substrate is PARTIAL-VIABLE
under fixed $r_d$ at $\lesssim 1.5\sigma$. It serves as the cross-bridge
anchor for the Phase~6e Sakharov 4-criterion test (§\ref{sec:phase6m-sakharov}).

%% ─── End Track C ───


\section{Sakharov criterion cross-bridge (final closure)}
\label{sec:phase6m-sakharov}

Phase~6m closes the Sakharov 4-criterion induced-gravity test introduced
in Phase~6e, providing the first systematic $\Lambda_J$ vs $\Lambda_{\rm HK}$
comparison on a common substrate in the literature. The test asks
whether the Jacobson cosmological constant $\Lambda_J$ equals the
Hawking-Killing $\Lambda_{\rm HK}$ value across substrate-level
predictions; the equality holds if and only if four substrate-level
conditions are satisfied:
\begin{enumerate}
\item Fermionic node structure (presence of low-energy fermion modes).
\item Universal coupling (substrate degrees of freedom couple universally
  to gravity).
\item Non-zero trace of identity ($\text{tr}\, I \neq 0$).
\item IR Lorentz recovery (the substrate's IR limit is Lorentz-invariant).
\end{enumerate}
The biconditional is captured in
\texttt{sakharov\_induced\_gravity\_criterion\_iff\_lambda\_j\_eq\_lambda\_hk}.

The substrate-level evaluations are:

\begin{itemize}
\item \texttt{volovik\_jannes\_he3a\_satisfies\_sakharov\_criterion}:
  Volovik-Jannes ${}^3\!\text{He-A}$ satisfies all four conditions
  (low-energy Weyl fermions; universal gravitational coupling via
  the tetrad; $\text{tr}\, I = 4$; IR Lorentz recovery). The
  substrate-derived $\Lambda$ value is
  $\Lambda_J = \Delta_0^4 / (6\pi^2 \hbar^3)$ at
  $\Delta_0 \approx \SI{1.6}{\milli\kelvin}$, with $\zeta = 2$.
\item \texttt{finazzi\_liberati\_sindoni\_bec\_violates\_universal\_coupling}:
  Finazzi-Liberati-Sindoni acoustic BEC violates condition~(ii) (universal
  coupling) --- the BEC's superfluid phonon mode does not couple
  universally to the substrate's effective gravitational sector at the
  required order, producing $\Lambda_J \neq \Lambda_{\rm HK}$.
\item \texttt{volovik\_jannes\_substrate\_anchors\_sakharov\_criterion}:
  the Volovik-Jannes substrate is the natural anchor for the Sakharov
  criterion; alternative substrates require explicit verification of
  all four conditions.
\end{itemize}

The substrate-class assignment + unimodular $\Lambda_{\rm HK}$ escape
route admits 5/6 Phase~6m candidates (NOT M8 KSS). The unimodular
reformulation
(\texttt{unimodular\_reformulation\_admits\_track\_c\_survivors}) provides
the secondary admissibility route. Track~C M8 (KSS) does not pass
through unimodular reformulation due to the Lorentz-violating
substrate; it remains conditional on path-(a).

\subsection*{§11.5 Sakharov $\Lambda_J$ vs $\Lambda_{\rm HK}$ (embedded short note)}
\label{subsec:sakharov-short-note}

The Sakharov 4-criterion test (§\ref{sec:phase6m-sakharov}) is the first
systematic $\Lambda$-comparison on a common substrate in the literature, and
constitutes a publishable observation in its own right. The full
short-note outline lives at
\texttt{temporary/working-docs/short\_note\_sakharov\_lambda\_j\_vs\_lambda\_hk\_outline.md}
(${\sim}\SI{20}{\kilo\byte}$, 32~references). It is embedded here per
Pipeline Invariant~\#14 default rather than spawning a stand-alone
CQG/PRD letter. The narrative covers: (i)~the four substrate-level
conditions and their physical content; (ii)~explicit numerics for
${}^3\!\text{He-A}$ ($\Lambda_J = \Delta_0^4 / (6\pi^2 \hbar^3)$,
$\Delta_0 \approx \SI{1.6}{\milli\kelvin}$, $\text{tr}\, I = 4$,
$\zeta = 2$); (iii)~the BEC counterexample's failure mode at
condition~(ii); (iv)~the unimodular reformulation as a $\Lambda_{\rm HK}$
escape route admitting 5/6 substrates; (v)~implications for the wider
emergent-gravity-DE program. User authorization for a 14th bundle
target would elevate this to a stand-alone Letter.


\section{Unified 7-class GD taxonomy and 3-tier applicability gradient}
\label{sec:phase6m-taxonomy}

The Phase~6m three-track closure is consolidated under a unified 7-class
GD taxonomy in \texttt{DarkSectorClassificationExtension.lean}. Each
mechanism is assigned to one of seven Gibbons--Duhem applicability classes:

\begin{itemize}
\item \emph{Class~0} (Track~A uniform): all causal-set DE mechanisms
  share Class~0 membership (substrate-d.o.f.\ outside GD domain;
  \texttt{track\_a\_uniform\_class\_0}; the membership is robust under
  admissible prescriptions per
  \texttt{class\_0\_membership\_robust\_under\_admissible\_prescriptions}).
\item \emph{Class~(a)} (Track~C M8 unique): KSS-style Lorentz-violating
  with anisotropic $c_s^2$ structure (\texttt{kss\_unique\_class\_a};
  uniqueness within Phase~6m).
\item \emph{Classes~(b), (c), (d)}: Track~B mechanism decomposition
  (\texttt{track\_b\_class\_decomposition}). GD binds only for
  Class~(c) and Class~(d) (\texttt{gd\_binds\_only\_for\_class\_c\_and\_class\_d}).
\item \emph{Classes~(e), (f), (g)}: Track~C survivor decomposition
  (\texttt{track\_c\_highest\_survival}); the 7-class structure
  organises all five+ CLEARED-R5 candidates by their substrate
  architecture.
\end{itemize}

The taxonomy supports a 3-tier applicability gradient
(\texttt{phase6m\_class\_admits\_tier\_assignment}):

\begin{itemize}
\item \emph{Tier 1 (GD applies, NO-GO predicted)}: Classes~(c), (d) ---
  Track~B's eight unanimous NO-GOs.
\item \emph{Tier 2 (GD silent, alternative analytic toolkit needed)}:
  Class~0 (causal-set DE), Class~(a) (KSS) --- the substrate
  d.o.f.\ are outside GD's Lagrangian-scalar domain.
\item \emph{Tier 3 (GD silent, geometric / cosmographic
  applicability)}: Classes~(e), (f), (g) --- Track~C's five+ CLEARED-R5
  survivors, where the analytic verdict comes from cosmographic
  consistency rather than thermodynamic consistency.
\end{itemize}

The 7-class taxonomy is the first unified organisational framework for
emergent-gravity DE mechanisms in the literature. It sets up the
flagship F~§8 dark-sector summary (which cites the 7-class structure as
the substrate program's dark-sector classification framework) and the
program's response to subsequent post-Phase-6m mechanism additions
(any new mechanism enters via class-assignment + tier-assignment, with
the orthogonality test against existing Phase~6m candidates following
mechanically).

The cross-bridge \texttt{class\_c\_admits\_orthogonality\_failure\_witness}
provides the existential witness that Class~(c) admits a substrate
configuration where the orthogonality principle fails --- the mechanism
that DEC instantiates in the Track~B closure. This is a Derived
content theorem (the witness is constructed explicitly, not
\texttt{decide}-on-hardcoded-flags); it is the rigorous companion to
DEC's failure-of-\texttt{IsViable} via the factor-\#1 chain.

The Phase~6m mechanism count is
\texttt{phase6m\_mechanism\_count}: 4 (Track~A) + 8 (Track~B) +
9 (Track~C) = 21 mechanisms evaluated, with the 8/8 unanimous
NO-GO closure on Track~B as the headline structural result.


\section{Discussion and outlook}
\label{sec:discussion}

The bundle's substantive contribution is the Phase~6m three-track
closure (§\ref{sec:phase6m-track-a}--\S\ref{sec:phase6m-taxonomy}),
which delivers (a)~Track~A's NO-GO+3-caveats verdict on causal-set DE,
(b)~Track~B's first-complete-mechanism-family unanimous NO-GO closure on
entropic-gravity DE, and (c)~Track~C's highest-survival verdict with
$f(R)$ Exp+ArcTanh as the strongest CLEARED-R5 of any track. The
Phase~6e Sakharov-criterion cross-bridge (§\ref{sec:phase6m-sakharov})
provides the program's first systematic $\Lambda_J$ vs $\Lambda_{\rm HK}$
test on a common substrate, falsifying acoustic-BEC analog gravity at
condition~(ii). The unified 7-class GD taxonomy
(§\ref{sec:phase6m-taxonomy}) organises all $21$ mechanisms under a
3-tier applicability gradient and sets up the response to subsequent
mechanism additions.

The bundle's structural caveats (Track~A GD-inapplicability under all
admissible prescriptions; Barrow-bound prescription dependence; BDG
$\sigma_\Lambda$ first-principles decomposition) are publishable
short-note opportunities embedded as D5~§10.5 + §11.5 per Pipeline
Invariant~\#14 default; user authorization would elevate them to
stand-alone CQG/JCAP and CQG/PRD Letters respectively.

DR3 + Roman-(${\sim}2030$) cross-track forecast:
pure-$\Lambda$ floating-$r_d$ disfavour reaches $\gtrsim 5\sigma$ at DR3
and ${>}7\sigma$ at Roman; fixed-$r_d$ stays $<3\sigma$ throughout
(fixed-vs-floating $r_d$ split = single empirical decider by 2030);
$f(R)$ Exp/ArcTanh hold at $\gtrsim 5\sigma$ preference; M4 Pure Lovelock
remains at the $1\sigma$-box edge. The closure programme thus has a
2-3-year quantitative outlook even for the conditional CLEARED-R5
candidates.

The bundle ships at $\totaltheorems$ Lean theorems / $\leanmodules$
modules / zero active sorry / 1 tracked axiom
(\texttt{gapped\_interface\_axiom} from Phase~5s, unrelated to the
dark sector). The Phase~6m closure adds $50$ substantive theorems
across 4 modules. The substantive count $\substantivetheorems$
separates load-bearing from $\placeholdertheorems$ documentation-marker
content.

\begin{thebibliography}{99}

%% ── paper17 sources (DM/DE classification + $\mathbb{Z}_{16}$) ──

\bibitem{Berezhiani2015}
L.~Berezhiani and J.~Khoury, ``Theory of dark matter superfluidity,''
Phys.\ Rev.\ D \textbf{92}, 103510 (2015), arXiv:1507.01019.

\bibitem{BK2025}
L.~Berezhiani, G.~Cintia, V.~De~Luca, and J.~Khoury,
``Superfluid Dark Matter,''
Physics Reports (2025), arXiv:2505.23900.

\bibitem{Volovik2006}
G.~E.~Volovik, ``Cosmological constant and vacuum energy,''
JETP Lett.\ \textbf{82}, 319 (2005), arXiv:gr-qc/0604062.

\bibitem{KV2008}
F.~R.~Klinkhamer and G.~E.~Volovik, ``Self-tuning vacuum variable and
cosmological constant,''
Phys.\ Rev.\ D \textbf{77}, 085015 (2008), arXiv:0711.3170.

\bibitem{ADW2019}
A.~A.~Abrikosov~Jr.\ and F.~R.~Klinkhamer,
``Tetrad condensation and emergent gravity,''
Phys.\ Rev.\ D \textbf{99}, 105009 (2019).

\bibitem{Diakonov2011}
D.~Diakonov,
``Towards lattice-regularized Quantum Gravity,''
arXiv:1109.0091 (2011).

\bibitem{Wan2019}
Z.~Wan, J.~Wang, and Y.~Zheng,
``Quantum 4d Yang--Mills theory and time-reversal symmetric 5d higher-gauge topological field theory,''
Phys.\ Rev.\ D \textbf{100}, 085012 (2019), arXiv:1904.00994.

\bibitem{DESI2024}
DESI Collaboration, ``DESI 2024 VI: cosmological constraints from the
measurements of BAO,'' arXiv:2404.03002 (2024).

\bibitem{DESI2025}
DESI Collaboration, ``DESI DR2 BAO: cosmological constraints,''
arXiv:2503.14738 (2025).

%% ── paper29 BBN sources ──

\bibitem{Pretko2017}
M.~Pretko, ``Subdimensional particle structure of higher rank
$U(1)$ spin liquids,'' Phys.\ Rev.\ B \textbf{95}, 115139 (2017),
arXiv:1604.05329.

\bibitem{Glodkowski2024}
A.~G\l{}\'odkowski, F.~Pe\~na-Ben\'itez, and P.~Sur\'owka,
``Dissipative fracton superfluids,'' arXiv:2401.01877 (2024).

%% ── paper32 Strong-CP / topological-DE sources ──

\bibitem{VanWaerbeke2025}
L.~Van~Waerbeke and A.~R.~Zhitnitsky,
``DESI results and Dark Energy from QCD topological sectors,''
arXiv:2506.14182 (2025).

%% ── Phase 6m three-track sources ──

\bibitem{PlazaKraiselburd2025fR}
M.~Plaza and L.~Kraiselburd,
``$f(R)$ gravity DESI DR2 + PPS + CC: very strong Bayesian preference,''
Phys.\ Rev.\ D \textbf{112}, 023554 (2025), arXiv:2504.05432.

\bibitem{TyagiHaridasuBasak2025}
S.~Tyagi, B.~S.~Haridasu, and S.~Basak,
``Bayesian model comparison of holographic dark-energy models with
DESI DR2,'' arXiv:2504.11308 (2025).

\bibitem{LucianoPaliathanasisSaridakis2506}
G.~G.~Luciano, A.~Paliathanasis, and E.~N.~Saridakis,
``Same entropies in holographic dark energy: Bayesian comparison with
DESI DR2,'' arXiv:2506.03019 (2025).

\bibitem{YoonGuha2023}
J.~Yoon and B.~Guha,
``Stability of de Sitter attractors in entropic FLRW cosmology with
matter and radiation,'' arXiv:2304.07301.

\bibitem{Vergeles2025}
S.~N.~Vergeles, ``Unitarity of 4D Lattice Theory of Gravity,''
Phys.\ Rev.\ D \textbf{112}, 054509 (2025), arXiv:2506.00036.

%% ── Phase 6m Lean infrastructure ──

\bibitem{Lean4}
L.~de~Moura and S.~Ullrich, ``The Lean 4 theorem prover and programming
language,'' CADE \textbf{28}, 625 (2021).

\bibitem{Mathlib}
The mathlib Community,
``The Lean mathematical library,''
CPP \textbf{2020}, 367 (2020), arXiv:1910.09336.

\end{thebibliography}

%% BUNDLE_APPEND_INSERT_HERE — bundle_append.py inserts new sections at
%% this marker by default. To insert at a specific section, edit the
%% body manually after the script runs.

\end{document}
